\begin{derivation}[从\eqnrefs{eq:qcd-gauge-field-transformation-dp10}到\eqnrefs{eq:qcd-ghost-expanded-dp68}]
取固定的 Lie 代数值函数 $\theta(x)=\theta^a(x)T^a$，并沿同一条规范轨道定义一参数曲线
\begin{equation*}
 U_\tau(x)=\exp[-\ii\tau\theta(x)],
 \qquad
 \mathscr C_\mu(\tau)
 =U_\tau\mathscr C_\mu U_\tau^{-1}
 +\frac{\ii}{g_s}(\partial_\mu U_\tau)U_\tau^{-1}.
\end{equation*}
在 $\tau=0$ 处有
\begin{equation*}
 \left.\frac{\dd U_\tau}{\dd\tau}\right|_0=-\ii\theta,
 \qquad
 \left.\frac{\dd U_\tau^{-1}}{\dd\tau}\right|_0=+\ii\theta .
\end{equation*}
对有限规范变换律直接求导：
\begin{align*}
 \left.\frac{\dd\mathscr C_\mu(\tau)}{\dd\tau}\right|_0
 &=-\ii[\theta,\mathscr C_\mu]
 +\frac1{g_s}\partial_\mu\theta\\
 &=\frac1{g_s}(D_\mu\theta)^aT^a,
\end{align*}
其中按照当前结构常数约定
\begin{equation*}
 (D_\mu\theta)^a
 =\partial_\mu\theta^a-g_sf^{abc}\mathscr C_\mu^b\theta^c.
\end{equation*}
若规范轨道参数采用反向曲线 $U_{-\tau}$，上式整体变号；以下把正文采用的轨道方向固定为
\begin{equation*}
 \delta_\theta\mathscr C_\mu^a
 \equiv
 \left.\frac{\dd\mathscr C_\mu^a(-\tau)}{\dd\tau}\right|_0
 =-\frac1{g_s}(D_\mu\theta)^a.
\end{equation*}
这只是对规范轨道坐标方向的定义，此后所有 Jacobian 都使用同一方向。

对 Lorenz 型规范函数
\begin{equation*}
 G^a[\mathscr C](x)=\partial^\mu\mathscr C_\mu^a(x),
\end{equation*}
其轨道导数为
\begin{equation*}
 \delta_\theta G^a(x)
 =-\frac1{g_s}\partial_x^\mu D_\mu^{ab}(x)\theta^b(x).
\end{equation*}
因此泛函导数严格等于
\begin{equation*}
 g_s\frac{\delta G^a(x)}{\delta\theta^b(y)}
 =-\partial_x^\mu D_\mu^{ab}(x)\delta^{(4)}(x-y)
 \equiv\mathcal M_{\rm FP}^{ab}(x,y),
\end{equation*}
即正文\eqnrefs{eq:qcd-fp-operator-dp10}。

在每条规范轨道与规范切片只有一个横截点的局部坐标片内，普通多维换元公式推广为
\begin{equation*}
 1=\Delta_{\rm FP}[\mathscr C]
 \int\mathcal D\theta\,\delta[G(\mathscr C^\theta)],
 \qquad
 \Delta_{\rm FP}[\mathscr C]=\det\mathcal M_{\rm FP},
\end{equation*}
其中与 $\mathscr C$ 无关的常数行列式已吸收到路径积分整体归一化。Grassmann Gaussian 恒等式给
\begin{equation*}
 \int\mathcal D\bar c\,\mathcal Dc\,
 \exp\!\left[
 \ii\int\dd^4x\,\bar c^a\mathcal M_{\rm FP}^{ab}c^b
 \right]
 =\mathcal N_{\rm gh}\det\mathcal M_{\rm FP},
\end{equation*}
其中 $\mathcal N_{\rm gh}$ 与背景规范场无关。把零背景值相除即可消去它：
\begin{equation*}
 \frac{\Delta_{\rm FP}[\mathscr C]}{\Delta_{\rm FP}[0]}
 =
 \frac{\displaystyle\int\mathcal D\bar c\,\mathcal Dc\,
 \ee^{\ii\int\dd^4x\,\bar c\mathcal M_{\rm FP}c}}
 {\displaystyle\int\mathcal D\bar c\,\mathcal Dc\,
 \ee^{\ii\int\dd^4x\,\bar c\mathcal M_{\rm FP}[0]c}}.
\end{equation*}
为了得到一般协变 $R_\xi$ 规范，引入辅助函数 $\omega^a(x)$ 以及归一化常数 $\mathcal N_\xi$，使
\begin{equation*}
 1=\mathcal N_\xi
 \int\mathcal D\omega\,
 \exp\!\left[-\frac{\ii}{2\xi}
 \int\dd^4x\,\omega^a\omega^a\right].
\end{equation*}
把 $\delta[G]$ 替换为 $\delta[G-\omega]$ 并先对 $\omega$ 积分，得到规范固定权重
\begin{equation*}
 \exp\!\left[-\frac{\ii}{2\xi}\int\dd^4x\,
 (\partial^\mu\mathscr C_\mu^a)^2\right],
\end{equation*}
对应
\begin{equation*}
 \mathcal L_{\rm gf}
 =-\frac{\hbar c}{2\xi}(\partial^\mu\mathscr C_\mu^a)^2.
\end{equation*}
与 Grassmann 行列式合并即得
\begin{equation*}
 \mathcal L_{\rm gf+gh}
 =-\frac{\hbar c}{2\xi}(\partial\!\cdot\!\mathscr C^a)^2
 +\hbar c\,\bar c^a(-\partial^\mu D_\mu^{ab})c^b.
\end{equation*}

最后展开伴随协变导数：
\begin{align*}
 \mathcal L_{\rm gh}
 &=\hbar c\,\bar c^a
 \left[-\partial^\mu
 \left(\partial_\mu c^a-g_sf^{abc}\mathscr C_\mu^bc^c\right)\right]\\
 &=-\hbar c\,\bar c^a\partial^2c^a
 +\hbar c g_sf^{abc}\bar c^a\partial^\mu(\mathscr C_\mu^bc^c).
\end{align*}
对第二项分部积分并取边界项为零：
\begin{equation*}
 \int\dd^4x\,\bar c^a\partial^\mu(\mathscr C_\mu^bc^c)
 =-\int\dd^4x\,(\partial^\mu\bar c^a)\mathscr C_\mu^bc^c.
\end{equation*}
于是
\begin{equation*}
 \boxed{
 \mathcal L_{\rm gh}
 =-\hbar c\,\bar c^a\partial^2c^a
 -\hbar c g_sf^{abc}(\partial^\mu\bar c^a)\mathscr C_\mu^bc^c,}
\end{equation*}
即正文\eqnrefs{eq:qcd-ghost-expanded-dp68}。第一项的逆核给鬼场传播子，第二项对 $\bar c,\mathscr C,c$ 的三阶泛函导数给鬼--反鬼--胶子顶角。
\end{derivation}
